\documentclass[a4paper,10pt]{article}
% Shared formatting for Anton Serdyuchenko's resumes.
% Compact single-column layout designed for readability and ATS-friendly text extraction.

\usepackage{fontspec}
\setmainfont{DejaVu Serif}
\usepackage{geometry}
\geometry{left=0.58in,right=0.58in,top=0.48in,bottom=0.48in}
\usepackage{enumitem}
\usepackage[hidelinks]{hyperref}
\usepackage{titlesec}
\usepackage{tabularx}
\usepackage{microtype}
\usepackage{xcolor}
\usepackage{ragged2e}

\pagestyle{empty}
\setlength{\parindent}{0pt}
\setlength{\tabcolsep}{0pt}
\urlstyle{same}
\raggedbottom
\setlength{\emergencystretch}{2em}

\titleformat{\section}
  {\vspace{-2pt}\large\bfseries\raggedright}
  {}{0em}{}
  [\titlerule\vspace{-4pt}]
\titlespacing*{\section}{0pt}{7pt}{4pt}

\newcommand{\resumeHeader}[3]{%
  \begin{center}
    {\LARGE\bfseries #1}\\[-1pt]
    {\normalsize\bfseries #2}\\[2pt]
    {\small #3}
  \end{center}
  \vspace{-7pt}
}

\newcommand{\resumeSubHeadingListStart}{%
  \begin{itemize}[leftmargin=0pt,label={},itemsep=5pt,topsep=0pt,parsep=0pt,partopsep=0pt]
}
\newcommand{\resumeSubHeadingListEnd}{\end{itemize}}

\newcommand{\resumeSubheading}[4]{%
  \item
  \begin{tabularx}{\textwidth}{X >{\raggedleft\arraybackslash}p{0.22\textwidth}}
    \textbf{#1} & {\small #2} \\
    \textit{#3} & \textit{\small #4}
  \end{tabularx}
  \vspace{-2pt}
}

\newcommand{\resumeProjectHeading}[3]{%
  \item
  \begin{tabularx}{\textwidth}{X r}
    \textbf{#1} \hspace{0.4em}{\small #2} & {\small #3}
  \end{tabularx}
  \vspace{-2pt}
}

\newcommand{\resumeItemListStart}{%
  \begin{itemize}[leftmargin=15pt,label=\textbullet,itemsep=1.5pt,topsep=2pt,parsep=0pt,partopsep=0pt]
}
\newcommand{\resumeItemListEnd}{\end{itemize}\vspace{-2pt}}
\newcommand{\resumeItem}[1]{\item {\small #1}}

\newcommand{\resumeLine}[2]{%
  \textbf{#1}: {\small #2}\\[-1pt]
}

\newcommand{\sep}{\hspace{0.45em}\textbar\hspace{0.45em}}


\begin{document}

\resumeHeader
  {Антон Сердюченко}
  {Java-разработчик | DevOps / AI Infrastructure}
  {\href{mailto:anton415@gmail.com}{anton415@gmail.com} \sep
   \href{https://www.linkedin.com/in/antonserdyuchenko}{LinkedIn} \sep
   \href{https://github.com/anton415}{GitHub} \sep
   \href{https://serdyuchenko.com}{serdyuchenko.com}}

\section{О СЕБЕ}
\small
Java-разработчик с 9-летним опытом разработки и сопровождения корпоративного ПО в Банке России. Основная работа: разработка на Java и Vaadin, исправление бизнес-логики, рефакторинг legacy-кода, JUnit-тесты, code review, SQL и диагностика поведения приложения. Параллельно развиваю практические навыки DevOps и AgentOps в open-source проекте \textit{finance-lab}: CI, автоматизация жизненного цикла PR, code review с coding agents, evals и human approval gates.

\section{ОПЫТ РАБОТЫ}
\resumeSubHeadingListStart
  \resumeSubheading
    {Банк России}{10.2017 -- настоящее время}
    {Ведущий инженер (Java-разработчик)}{Россия}
    \resumeItemListStart
      \resumeItem{Разрабатываю и поддерживаю функциональность долгоживущего корпоративного приложения, в основном на Java и Vaadin, включая бизнес-логику и UI.}
      \resumeItem{Реализовал сквозной рефакторинг одного из ключевых бизнес-процессов: изменения доменной логики, UI, сервисов и workflow orchestration с переходом затронутой функциональности на обновленную доменную модель; проверил изменения JUnit- и ручными тестами.}
      \resumeItem{Разработал переиспользуемые Vaadin-компоненты для повторяющихся UI-паттернов в нескольких формах приложения.}
      \resumeItem{Исправляю дефекты бизнес-логики и рефакторю legacy-код; использую SQL и диагностику приложения для анализа и проверки поведения.}
      \resumeItem{Провожу code review изменений коллег и пишу или поддерживаю JUnit-тесты в рамках регулярной разработки и проверки изменений.}
    \resumeItemListEnd
\resumeSubHeadingListEnd

\section{ПРОЕКТНАЯ РАБОТА}
\resumeSubHeadingListStart
  \resumeProjectHeading
    {finance-lab -- open-source лаборатория AgentOps}
    {\href{https://github.com/anton415/finance-lab}{GitHub}}
    {2026 -- настоящее время}
    \resumeItemListStart
      \resumeItem{Создал open-source приложение для личных финансов и инженерную лабораторию для экспериментов с надежной AI-agent-assisted разработкой.}
      \resumeItem{Выстроил процесс, в котором coding agents реализуют и проверяют изменения, а человек сохраняет контроль над acceptance criteria, review, merge и внешними действиями.}
      \resumeItem{В проекте задокументирован evaluation framework из 12 candidate tasks; для одной задачи проведена ручная калибровка grader с positive/negative controls.}
      \resumeItem{Использую GitHub Actions, инструкции для агентов на уровне репозитория и Python-автоматизацию с разделением dry-run/live для наблюдаемого и контролируемого workflow.}
    \resumeItemListEnd
\resumeSubHeadingListEnd

\section{НАВЫКИ}
\resumeLine{Разработка}{Java, Vaadin, SQL, JUnit, Git/GitLab, code review, debugging, legacy refactoring}
\resumeLine{Backend / Workflow}{PostgreSQL, Camunda/BPMN, Mockito, Maven}
\resumeLine{Delivery / инфраструктура}{GitHub Actions, CI/CD, Linux, базовая Python-автоматизация}
\resumeLine{AI / Agent Engineering}{coding-agent workflows, eval design, human approval gates, AgentOps experiments}

\section{ОБРАЗОВАНИЕ}
\resumeSubHeadingListStart
  \resumeSubheading
    {Московский авиационный институт (МАИ)}{2022}
    {Магистратура: менеджмент, программа «Экономическая безопасность», факультет ИТ}{Москва}
\resumeSubHeadingListEnd

\section{ДОПОЛНИТЕЛЬНОЕ ОБУЧЕНИЕ}
\small
Camunda Academy -- BPMN Overview \sep
Anthropic -- AI Fluency Framework \& Foundations \sep
Anthropic -- Claude Code in Action

\section{ЯЗЫКИ}
\resumeLine{Русский}{родной}
\resumeLine{Английский}{B2}

\end{document}
